% =============================================================================
%  When does FedProx help? — synthesis subsection
%
%  Intended as a subsection of \section{Discussion} or a standalone synthesis
%  near the start of \section{Conclusion}. Reads as a structured argument:
%  one row per empirical condition tested, with the observed Δ, the
%  inferential verdict, and the mechanistic interpretation.
%
%  Placeholders to fill once the canonical analysis pipeline has run
%  on the full HPC dataset:
%    \TODOhpc{value} — replace with the post-analysis number
%    \TODOhpc{p-value} — paired Wilcoxon p
%    \TODOhpc{interp} — narrative interpretation
%
%  Cross-references assume these labels exist:
%    sec:stat-het          — the statistical-heterogeneity experiments section
%    sec:sys-het           — the system-heterogeneity experiments section
%    sec:results-statistical, sec:results-systemhet — §Results subsections
%    sec:hyp-specialist    — the specialist pre-registration
%    fig:dose-response     — the measured-heterogeneity dose-response figure
% =============================================================================

\subsection{When does FedProx help, and when doesn't it?}
\label{sec:when-fedprox}

The empirical question this dissertation set out to answer is not
\emph{whether} FedProx outperforms FedAvg in some abstract sense — the
asymptotic theory in \citet{li2020fedprox} already provides one answer
to that — but rather \emph{under what specific conditions} the proximal
anchor produces a measurable improvement on a non-trivial medical-imaging
classification task. This subsection synthesises the empirical evidence
from the experimental matrix (\S\ref{sec:stat-het}, \S\ref{sec:sys-het})
into a per-condition verdict.

\subsubsection{Decision table: per-condition verdicts}
\label{sec:when-fedprox-table}

Table~\ref{tab:when-fedprox} reports the test-set macro-F1
within-pair difference $\Delta = \mathrm{FedProx} - \mathrm{FedAvg}$
under each experimental condition. The verdict column distinguishes:

\begin{itemize}
  \item \textbf{HELPS}: $\Delta > 0$ at $\alpha = 0.05$ (paired Wilcoxon
        two-sided); effect size $|r_{\mathrm{rb}}| \geq 0.5$
        ("very large" by Cohen's conventions for rank-biserial).
  \item \textbf{NULL}: the data do not warrant rejecting
        $H_0: \Delta = 0$; the exact Walsh-average 95\% confidence
        interval for the Hodges--Lehmann estimate of
        $\mathrm{median}(\Delta)$ is contained in
        $[-\mathrm{SES}, +\mathrm{SES}]$ with
        $\mathrm{SES} = 0.01$ macro-F1 the pre-registered smallest
        effect of interest.
  \item \textbf{INCONCLUSIVE}: the data do not reject $H_0$, but the
        CI extends beyond $\pm\mathrm{SES}$; underpowered to
        distinguish "no effect" from "small effect".
  \item \textbf{HURTS}: $\Delta < 0$ with paired Wilcoxon
        two-sided $p < 0.05$.
\end{itemize}

\begin{table}[ht]
\centering
\small
\caption{When FedProx beats FedAvg on DermaMNIST at $E=20$, $R=150$, $\mu=0.01$,
$n = 10$ paired seeds, two-sided paired Wilcoxon, smallest effect of interest
$\mathrm{SES} = 0.01$. Verdict column codes: \textbf{HELPS} ($\Delta > 0$,
$p < 0.05$); \textbf{NULL} (CI for $\mathrm{HL}(\Delta) \subset [-0.01, +0.01]$);
\textbf{INCONCLUSIVE} (no rejection, but CI exceeds SES);
\textbf{HURTS} ($\Delta < 0$, $p < 0.05$).}
\label{tab:when-fedprox}
\begin{tabular}{l c c c l}
\toprule
Condition & $\mathrm{HL}(\Delta)$ & 95\% CI & $p_{\mathrm{Wilcoxon}}$ & Verdict \\
\midrule
\multicolumn{5}{l}{\emph{Statistical heterogeneity}} \\
balanced\_paired\_7\_clients (headline)
   & \TODOhpc{HL_paired}     & \TODOhpc{CI_paired}     & \TODOhpc{p_paired}     & \TODOhpc{HELPS} \\
specialist\_7\_clients (1-owner/class)
   & \TODOhpc{HL_spec}       & \TODOhpc{CI_spec}       & \TODOhpc{p_spec}       & \TODOhpc{NULL/INCONCL} \\
dirichlet\_alpha01\_7\_clients (random)
   & \TODOhpc{HL_dir}        & \TODOhpc{CI_dir}        & \TODOhpc{p_dir}        & \TODOhpc{NULL/INCONCL} \\
iid\_7\_clients (no skew)
   & \TODOhpc{HL_iid}        & \TODOhpc{CI_iid}        & \TODOhpc{p_iid}        & \TODOhpc{NULL} \\
\midrule
\multicolumn{5}{l}{\emph{System heterogeneity (balanced\_paired partition fixed)}} \\
C0 uniform ($E_i = 20$ all clients)
   & \TODOhpc{HL_C0}         & \TODOhpc{CI_C0}         & \TODOhpc{p_C0}         & \TODOhpc{verdict} \\
C1 fixed stragglers (C5,C6 at $E=5$)
   & \TODOhpc{HL_C1}         & \TODOhpc{CI_C1}         & \TODOhpc{p_C1}         & \TODOhpc{verdict} \\
C2 random stragglers ($E_i \sim \mathrm{Unif}\{1..19\}$)
   & \TODOhpc{HL_C2}         & \TODOhpc{CI_C2}         & \TODOhpc{p_C2}         & \TODOhpc{verdict} \\
\midrule
\multicolumn{5}{l}{\emph{Sensitivity sweeps (3 seeds per cell)}} \\
$\mu$ varying \{0.001, 0.01, 0.1, 0.5, 1.0\}
   & \multicolumn{4}{l}{see Figure~\TODOhpc{fig:mu-sweep}; peak at $\mu^* = \TODOhpc{mustar}$} \\
$E$ varying \{1, 5, 10, 20, 40\}
   & \multicolumn{4}{l}{see Figure~\TODOhpc{fig:e-sweep}; slope $\partial \Delta / \partial \log E = \TODOhpc{slope}$} \\
\bottomrule
\end{tabular}
\end{table}

\subsubsection{Three regimes characterised}
\label{sec:when-fedprox-regimes}

Reading Table~\ref{tab:when-fedprox} from top to bottom, the
empirical evidence places FedProx into one of three regimes
relative to FedAvg:

\paragraph{Regime A — FedProx helps decisively.}
\TODOhpc{Decide which condition(s) qualify based on the actual
data; pre-registered candidate is balanced\_paired only.}
Conditions in this regime share a specific structural feature:
\emph{at least two clients hold overlapping minority-class data
and therefore drift apart during local SGD}. The proximal anchor
provides a per-round consensus-finding mechanism at the
aggregation step that the unconstrained FedAvg averaging cannot
match. \TODOhpc{Cite the relevant per-class Holm-survivor — likely
melanoma — to anchor the mechanism story to a concrete observable.}

\paragraph{Regime B — FedProx is indistinguishable from FedAvg.}
\TODOhpc{Likely IID and specialist; possibly Dirichlet depending
on observed CI bounds.} In these conditions, either the local
objectives $F_i$ are too similar to the global $F$ for the
proximal anchor to have meaningful work to do (the IID case), or
the heterogeneity structure does not produce the paired
drift-apart pattern that the anchor reconciles (the specialist
and possibly Dirichlet cases). The proximal term reduces to a
constant penalty that effectively damps individual-client
gradient magnitudes by a small factor — equivalent in expectation
to a slight learning-rate reduction — which neither helps nor
harms on average.

\paragraph{Regime C — FedProx might hurt (boundary case).}
\TODOhpc{If any condition shows $\Delta < 0$ at $p < 0.05$,
discuss here. Pre-registered candidate based on partial data:
C1 fixed stragglers (where one outlier seed showed catastrophic
FedProx failure). If no condition meets HURTS criterion after
full data, state that explicitly: "No condition in our matrix
demonstrated a statistically significant FedProx disadvantage."}
The mechanism for a negative effect, if observed, would be
over-anchoring: when a single client holds the only samples of a
rare class, pulling its local update back toward a global
parameter vector that under-represents that class can actively
delay learning for that class.

\subsubsection{Implications for FL deployment in medical imaging}
\label{sec:when-fedprox-deployment}

For a practitioner deciding whether to deploy FedProx in a real
cross-silo federation, the empirical recommendation from this
thesis is:

\begin{itemize}
  \item \textbf{Adopt FedProx} when the federation contains
        multiple sites each contributing overlapping rare-class
        data (the most common cross-silo medical pattern: many
        hospitals each contributing a few rare-disease cases).
        This is the paired-co-training regime where the proximal
        anchor's gain is largest.
  \item \textbf{FedAvg is sufficient} when the federation
        approximates IID (large geographically dispersed
        consumer-FL setting, well-mixed data) or when each
        site is a sole specialist holding a single rare class
        (no co-training to reconcile).
  \item \textbf{Tune $\mu$ on a small per-federation pilot}
        rather than assuming the default value transfers across
        deployments. Our $\mu$-sensitivity sweep
        (Figure~\TODOhpc{fig:mu-sweep}) shows that the
        FedProx advantage \TODOhpc{is/is-not} robust to
        $\mu$ choice over an order of magnitude.
  \item \textbf{Account for run-to-run variation.} The
        cross-runtime divergence observed between our
        pure-PyTorch reference loop and the Flower simulation
        runtime (\TODOhpc{Δ_PT = ..., Δ_Flower = ...}) suggests
        that the FedProx-vs-FedAvg gap at modest sample sizes
        is more runtime-sensitive than the prior literature has
        documented. Practitioners should validate any reported
        FedProx-vs-FedAvg comparison on the specific FL
        framework they intend to deploy.
\end{itemize}

These deployment-level recommendations follow directly from the
per-condition verdicts in Table~\ref{tab:when-fedprox}. They are
deliberately conservative: each prescription is grounded in a
specific empirical row of the table rather than extrapolated from
the headline result alone.

\subsubsection{Limits of the present characterisation}
\label{sec:when-fedprox-limits}

The "when does FedProx help?" answer above is bounded by what the
experimental matrix actually tested:

\begin{itemize}
  \item \emph{Single dataset, single architecture}: DermaMNIST
        $28\times28$ with the \texttt{DermMNISTCNN} 4-block
        convolutional model. Whether the same per-condition
        verdicts hold at larger resolutions, on different
        modalities, or with different architectures is an open
        question (see \S\ref{sec:future} for a 64$\times$64
        resolution sensitivity proposal).
  \item \emph{$K = 7$ clients, $C = 1.0$}: the cross-silo
        full-participation regime. Cross-device FL (large $K$,
        small $C$) is outside this thesis's scope; the
        canonical-comparator literature \citep{kairouz2021advances}
        characterises that regime separately.
  \item \emph{$n = 10$ paired seeds}: a meaningful but modest
        sample. Effect sizes near the smallest-effect-of-interest
        ($|\Delta| \approx 0.01$) are unlikely to be resolved at
        this $n$; the INCONCLUSIVE verdict rows in
        Table~\ref{tab:when-fedprox} are honest about that
        limitation rather than over-claiming "no effect".
  \item \emph{Two runtimes evaluated}: the pure-PyTorch reference
        loop (the source of the historical headline) and the
        Flower 1.x simulation. Other FL frameworks (NVFlare,
        FedML, FedScale) may produce different numerical
        trajectories; we have not tested them.
\end{itemize}

These boundaries are not weaknesses to apologise for — they are
the precise statement of what the four-partition seven-client
$E$=20-$R$=150 evaluation can and cannot say. A reviewer
asking "but does FedProx help in setting X?" can be answered
honestly: \emph{X is or is not in our experimental matrix;
if not, see the relevant future-work direction in
\S\ref{sec:future}}.
